\chapter{Thesis Statement and Research Questions}
\label{chap:thesis_statement}

\section{Thesis Statement}
\label{sec:thesis_statement}

This thesis argues that a \textit{hybrid AI framework} — combining symbolic AI (specifically ontology-based knowledge graphs using standards such as ESCO) with subsymbolic AI (data-driven forecasting and generative AI) — can serve as an effective foundation for an intelligent recommendation engine that improves resource demand forecasting and allocation decision-making in CAPEX portfolio management. The system aims to enable organizations to proactively identify skill bottlenecks and evaluate strategic workforce responses (make, buy, or upskill) before resource conflicts materialize.

\section{Research Questions}
\label{sec:research_questions}

\begin{description}
  \item[RQ1] How can a hybrid AI approach — integrating ontology-based knowledge representation with data-driven machine learning — be designed to forecast skill-specific resource demand across CAPEX project lifecycle phases?

  \item[RQ2] Which symbolic knowledge structures (e.g., skill ontologies such as ESCO, project taxonomy graphs) are most suitable for enriching subsymbolic demand forecasting models in the CAPEX resource management domain?

  \item[RQ3] How can detected bottlenecks be translated into actionable, constraint-aware allocation recommendations that account for organizational make/buy/upskill options?

  \item[RQ4] To what extent does the proposed hybrid AI recommendation engine outperform conventional planning approaches in terms of bottleneck anticipation lead time and resource utilization efficiency?
\end{description}
